% Глава 1: Введение / Chapter 1: Introduction

\chapter{Введение}
\label{ch:introduction}

\section{Актуальность темы исследования}
\label{sec:intro:relevance}

[TODO: Описать актуальность проблемы ввода и верификации данных в наукометрических системах]

Интерактивные наукометрические системы, такие как ИСТИНА МГУ~\cite{TODO:istina}, Web of Science~\cite{TODO:wos}, Scopus~\cite{TODO:scopus}, играют критическую роль в современной научной инфраструктуре. Эти системы обрабатывают миллионы библиографических записей, позволяя исследователям, администраторам и государственным органам оценивать научную продуктивность, выявлять тенденции и принимать обоснованные решения о финансировании.

Однако ввод и верификация больших объемов данных в таких системах остаются существенной проблемой по следующим причинам:

\begin{enumerate}
    \item \textbf{Многоязычность.} Авторы публикаций используют имена на разных языках (латиница, кириллица, иероглифы), что приводит к вариативности написания и неоднозначности интерпретации.

    \item \textbf{Неструктурированность метаданных.} Источники данных (Crossref~\cite{TODO:crossref}, ORCID~\cite{TODO:orcid}) часто содержат неполные или некорректные метаданные, требующие дополнительной верификации.

    \item \textbf{Масштаб данных.} Современные наукометрические системы обрабатывают сотни тысяч записей ежегодно, что делает ручную верификацию практически невозможной.

    \item \textbf{Культурная специфика.} Различные традиции именования (например, китайская, корейская, вьетнамская) требуют специализированных алгоритмов обработки.
\end{enumerate}

\section{Цель и задачи исследования}
\label{sec:intro:objectives}

\textbf{Цель исследования:}
Разработка методов и алгоритмов автоматизированного ввода и верификации больших данных в интерактивных наукометрических системах.

\textbf{Задачи исследования:}

\begin{enumerate}
    \item Провести анализ существующих подходов к автоматизации ввода и верификации данных в наукометрических системах.

    \item Разработать алгоритм автоматической верификации порядка имени и фамилии в китайских именах авторов.

    \item Реализовать мультимодульную архитектуру принятия решений с цепочками доказательств для повышения интерпретируемости результатов.

    \item Провести экспериментальную оценку точности и производительности разработанного метода на реальных данных.

    \item Внедрить разработанное решение в промышленную эксплуатацию системы ИСТИНА МГУ.
\end{enumerate}

\section{Научная новизна}
\label{sec:intro:novelty}

В диссертации получены следующие новые результаты:

\begin{enumerate}
    \item Предложен новый подход к автоматической верификации порядка имени и фамилии в китайских именах на основе мультимодульного алгоритма принятия решений.

    \item Разработана архитектура с цепочками доказательств, обеспечивающая полную интерпретируемость каждого решения системы.

    \item Создан механизм автоматической редакции персональных данных (PII redaction), соответствующий принципам GDPR и ФЗ-152.

    \item Экспериментально подтверждена эффективность предложенного метода на крупномасштабных реальных данных (N=301~446).
\end{enumerate}

\section{Практическая значимость}
\label{sec:intro:significance}

Результаты исследования внедрены в промышленную эксплуатацию системы ИСТИНА МГУ. Разработанный алгоритм обрабатывает более 10~000 китайских имен авторов в рамках регулярного импорта данных, обеспечивая точность [TODO: insert value]~\% и пропускную способность [TODO: insert value]~имен в секунду.

Кроме того, предложенная архитектура с цепочками доказательств может быть адаптирована для решения других задач верификации данных в интерактивных системах.

\section{Положения, выносимые на защиту}
\label{sec:intro:defense}

\begin{enumerate}
    \item Мультимодульный алгоритм верификации порядка имени и фамилии в китайских именах с использованием лингвистических, статистических и контекстных признаков.

    \item Архитектура принятия решений с цепочками доказательств, обеспечивающая полную интерпретируемость и аудируемость результатов.

    \item Механизм автоматической редакции персональных данных с использованием хеширования SHA-256 с солью.

    \item Результаты экспериментальной оценки точности и производительности на крупномасштабных реальных данных.
\end{enumerate}

\section{Апробация работы}
\label{sec:intro:approbation}

[TODO: Указать конференции, публикации, патенты]

Основные результаты работы были представлены на следующих конференциях и семинарах:

\begin{itemize}
    \item [TODO: Conference 1]
    \item [TODO: Conference 2]
    \item [TODO: Seminar/Workshop]
\end{itemize}

По теме диссертации опубликовано [TODO: N] работ, из них [TODO: M] — в журналах, рекомендованных ВАК.

\section{Структура диссертации}
\label{sec:intro:structure}

Диссертация состоит из введения, шести глав основного содержания, заключения и списка литературы.

\textbf{Глава 1 (Введение)} содержит обоснование актуальности темы, формулировку цели и задач, описание научной новизны и практической значимости работы.

\textbf{Глава 2 (Обзор литературы)} посвящена анализу существующих подходов к автоматизации ввода и верификации данных, обработке многоязычных имен и принятию решений в интерактивных системах.

\textbf{Глава 3 (Методология)} описывает предложенный мультимодульный алгоритм, архитектуру с цепочками доказательств и механизм редакции персональных данных.

\textbf{Глава 4 (Экспериментальные результаты)} представляет результаты экспериментальной оценки на крупномасштабных данных Crossref (N=301~446).

\textbf{Глава 5 (Интеграция с системой ИСТИНА)} описывает внедрение разработанного решения в промышленную эксплуатацию системы ИСТИНА МГУ.

\textbf{Глава 6 (Обсуждение)} содержит анализ полученных результатов, сравнение с существующими методами и обсуждение ограничений.

\textbf{Глава 7 (Заключение)} подводит итоги работы и намечает направления дальнейших исследований.
